\section{Sets express the boundary of a question}
A set is a collection whose membership matters. The set $A=\{1,3,5\}$ has three elements. Repeating a member does not create a new element: $\{1,1,3\}=\{1,3\}$. Order does not matter either. By contrast, an ordered pair $(1,3)$ is different from $(3,1)$. The distinction between a collection and an ordered arrangement is part of the mathematical model.

Write $x\in A$ when $x$ is a member and $x\notin A$ otherwise. The empty set $\varnothing$ has no elements. The notation $A\subseteq B$ says that every member of $A$ is a member of $B$. To prove this containment, take an arbitrary $x\in A$ and show $x\in B$. To prove equality of sets, prove containment in both directions.

The union $A\cup B$ contains elements in at least one set. The intersection $A\cap B$ contains elements in both. The difference $A\setminus B$ contains elements in $A$ but not in $B$. For $A=\{1,2,3\}$ and $B=\{3,4\}$, these are $\{1,2,3,4\}$, $\{3\}$, and $\{1,2\}$ respectively.

The notation $\{x\in A:P(x)\}$ selects the members of $A$ satisfying a property $P$. For example, $\{n\in\ZZ:n\text{ is even}\}$ is the set of even integers. Read the colon as ``such that.''

\subsection*{Why a membership proof proves an entire set equality}
The statements $1\in\{1,2\}$ and $\{1\}\subseteq\{1,2\}$ have different subjects. The first talks about one number. The second talks about a set whose members must all occur in another set. In particular, ``$x$ is a member'' and ``the one-element set $\{x\}$ is contained'' are related but differently written claims.

Two sets are equal precisely when they have the same members. To prove $D=E$, it is therefore enough to show that anything in $D$ is in $E$, and anything in $E$ is in $D$. Start the first direction with an arbitrary $x\in D$; do not assume at the outset that it also belongs to $E$. Producing that second membership is the work of the argument.

If the starting set is empty, there is no member to defeat the proposed containment. This is why $\varnothing\subseteq E$ for every set $E$. The sentence ``take an arbitrary member'' is conditional on a member being supplied; it does not assert that the set is nonempty. This is the set version of vacuous truth.

For a concrete preview of the next identity, use $A=\{1,2,3\}$, $B=\{2,4\}$, and $C=\{3,4\}$. First $B\cup C=\{2,3,4\}$; retaining members also in $A$ leaves $\{2,3\}$. On the other side, $A\cap B=\{2\}$ and $A\cap C=\{3\}$; their union is again $\{2,3\}$. The proof below explains why this agreement persists for arbitrary sets.

\par\noindent\begin{minipage}{\linewidth}
\begin{proposition}
For any sets $A,B,C$, we have $A\cap(B\cup C)=(A\cap B)\cup(A\cap C)$.
\end{proposition}
\end{minipage}\par

\begin{proof}
Suppose $x\in A\cap(B\cup C)$. Then $x\in A$ and $x$ belongs to $B$ or to $C$. In the first case $x\in A\cap B$; in the second $x\in A\cap C$. In either case $x$ belongs to their union. This proves one containment.

Conversely, suppose $x\in(A\cap B)\cup(A\cap C)$. In the first case $x\in A$ and $x\in B$; in the second $x\in A$ and $x\in C$. Both cases imply $x\in A$ and $x\in B\cup C$. Thus $x\in A\cap(B\cup C)$, proving the reverse containment.
\end{proof}
The proof did not depend on what the elements were. Its content was the pattern of membership. This is a first example of a structure that survives a change of subject matter.

\subsection*{Counting members and forming pairs}
For a finite set $S$, the notation $|S|$ means the number of its members. For example, if $S=\{2,4,6\}$, then $|S|=3$, not twelve. We count members, not their sum. The empty set has size zero. The word \emph{cardinality} is another name for this size. Bars around a real number mean its absolute value; bars around a finite set mean its cardinality. The kind of object inside the bars tells us which meaning is being used.

Two sets are \emph{disjoint} if they have no member in common. Counting a disjoint union adds the two sizes. If the sets overlap, simply adding their sizes counts the common members twice. For example, $|\{1,2\}|+|\{2,3\}|=4$ but $|\{1,2,3\}|=3$.

The \emph{Cartesian product} $A\times B$ is the set of ordered pairs $(a,b)$ with $a\in A$ and $b\in B$. If $A=\{1,2\}$ and $B=\{p,q\}$, it contains the four pairs $(1,p),(1,q),(2,p),(2,q)$. In general, each of the $|A|$ choices of first entry can be combined with each of the $|B|$ choices of second entry, so a finite product has $|A\times B|=|A|\cdot|B|$ members. Despite the multiplication sign, $A\times B$ does not multiply any numbers in $A$ and $B$; it forms pairs of objects, and only the \emph{count} of those pairs is a product. A rule whose input consists of two points of a set $X$, such as the distance between two points, can be described as a function on $X\times X$.
